\section{A small common library for represented mathematics}
\label{sec:analytic-tools}

\subsection{Removing a limiting phrase from a rational conclusion}

\begin{lemma}[Rational separation]\label{lem:rational-separation}
Let \(a,b\) be rationals, and let \(D\ge1\) be a common positive
denominator.  If \(a\le b+2^{-p}\) for
\(p=\operatorname{bl}(D)+1\), where \(\operatorname{bl}(D)\) is binary length,
then \(a\le b\).
Consequently, a proof of \(a\le b+2^{-p}\) at every precision implies
\(a\le b\) by one explicitly selected finite instance.
\end{lemma}

\begin{proof}
The comparison of \(a\) and \(b\) is decidable.  If \(a>b\), their
common-denominator difference is a positive integer divided by \(D\),
and is at least \(1/D\).  But \(2^{\operatorname{bl}(D)}>D\), so the specified error is
strictly smaller than \(1/D\), a contradiction.  This uses decidable
rational cases, not excluded middle for an infinite name.
\end{proof}

The same argument turns \(|a|\le2^{-p}\) for all \(p\) into an exact
rational equality.  A related integer argument turns an integer error
of absolute value smaller than one into zero.  For a represented real
there is no finite common denominator; its equality remains the family
of all finite precision statements.  This elementary distinction is what
justifies the rational limit-elimination steps used throughout the report.

\subsection{Uniform finite precision propagation}

The following lemma removes repeated informal instructions to
``take the approximation sufficiently accurate.''  It is a theorem about
rational arithmetic tables; represented values are handled only after
this finite theorem is available.

\begin{lemma}[A precision budget for a finite arithmetic table]
\label{core:gate}
Consider a directed arithmetic table of \(N\) gates, each an addition,
subtraction, or multiplication of previous entries.  Its inputs and
constants have supplied rational magnitude bounds.  Assign a bound to
each gate by adding its input bounds for addition or subtraction, and
multiplying them for multiplication.  Let the integer \(B\ge1\) dominate
all these recursively computed bounds.  Put \(L=2B+3\).

Suppose corresponding inputs differ by at most \(\delta\), and the second
evaluation introduces additional local rounding errors at most
\(\delta\) per gate.  If
\[
 \delta\le \frac{2^{-s}}{L^{N+1}},
\]
then the output error is at most \(2^{-s}\), and the error after the first
\(j\) gates is at most \(L^j\delta\).  All quantities in this prescription
are primitive recursive in the finite table and its bounds.
\end{lemma}

\begin{proof}
Let \(e_j\) be the largest error among inputs and the first \(j\) gates.
The input hypothesis gives \(e_0\le\delta\).  While \(e_j\le1\), a first
evaluation operand has magnitude at most \(B\), and its corresponding
second operand has magnitude at most \(B+1\).  For multiplication,
\[
 |ab-a'b'|
 \le |a|\,|b-b'|+|b'|\,|a-a'|
 \le(2B+1)e_j.
\]
Addition or subtraction gives at most \(2e_j\).  Adding the local
rounding error, and using \(\delta\le L^j\delta\), shows that the next
error is at most \(L^{j+1}\delta\).  The hypothesis makes this bound at
most one throughout the finite table, so the side condition is preserved.
At \(j=N\), it is at most \(2^{-s}/L\), which is smaller than the stated
output allowance.  The invariant is the conjunction of finitely many
rational bounds on the processed entries and their errors.  Its Boolean
fold is a quantifier-free formula.  Thus the displayed induction is
available in \ipra.
\end{proof}

For example, if \(b(t)\) is the binary length, the common input precision
\[
 p=s+b(L^{N+1})+2
\]
gives a dyadic error smaller than the required \(\delta\).
This is intentionally generous.  A symbolic error expression obtained
gate by gate often gives a much smaller precision.  Both schedules are
PR; no search for a successful precision is involved.  Bounds computed
from the syntax of a \emph{finite arithmetic table} do not evaluate an
arbitrary PR source program: producing that table has its own declared
PR construction or clock.

Guarded division can be added with a separate premise.  For real or complex
values in the usual absolute-value norm, if
\(|w|\ge\alpha>0\) and \(|w-w'|\le\eta\le\alpha/2\), then
\begin{equation}\label{eq:reciprocal-budget}
 |w'|\ge\alpha/2,\qquad
 \left|\frac1w-\frac1{w'}\right|\le\frac{2\eta}{\alpha^2}.
\end{equation}
The first inequality is the triangle inequality.  The second follows
from \((w'-w)/(ww')\).  Replace the local multiplier by a rational bound
on \(2/\alpha^2\), and preserve the denominator guard as part of the loop
invariant.  For complex computations one can check lower bounds for
squared Euclidean norms using rational arithmetic.  All denominator
guards must actually be supplied or proved.

\subsection{Some operations and their explicit resources}

For regular real names with stage error at most \(2^{-p}\), the following
constructions illustrate the common interface.  The error descriptions
refer to the usual interpretation; their proofs are the same finite
rational comparisons at two requested precisions.
\begin{center}\small
\begin{tabular}{@{}P{.18\linewidth}P{.25\linewidth}P{.48\linewidth}@{}}
\toprule
Operation & Input information & Construction and output error\\
\midrule
Addition & Two regular names & Add their stages \(k+2\); error at most
\(2^{-k-1}\).\\
Multiplication & \(|x|,|y|\le M\), integer \(M\ge1\) & Multiply their
stages \(p=k+2+b(2M+1)\).  The error is at most
\((2M+1)2^{-p}\le2^{-k-2}\).\\
Reciprocal & \(|x|\ge2^{-g}\) & Invert stage \(p=k+2g+4\).
It is nonzero, and \eqref{eq:reciprocal-budget} gives error at most
\(2^{-k-3}\).\\
Nonnegative square root & \(0\le x\le M\), integer \(M\ge1\) & Clip stage
\(2k+6\) to \([0,M]\); bisect its rational square root in \([0,M+1]\)
for \(k+3+b(M+1)\) steps.  Error is at most \(2^{-k-2}\).\\
Polynomial evaluation & Coefficient and argument bounds & Use the finite
table and Lemma~\ref{core:gate}; coefficients and degree are finite data.\\
\bottomrule
\end{tabular}
\end{center}

Here is the only nonrational estimate needed for the square-root entry.
If \(u,v\ge0\), then
\((\sqrt u-\sqrt v)^2\le |u-v|\), since, after ordering the roots,
\((a-b)^2\le(a-b)(a+b)\).  It can equally be read as an inequality for
two rational root brackets.  Clipping cannot increase error from a value
in \([0,M]\), so the input error contributes at most \(2^{-k-3}\) to
the root error.  Each bisection halves the rational bracket, and the
stated number of steps makes its width at most \(2^{-k-3}\).  Rational
comparison of the squared midpoint decides each step.  This also shows
why a square root at zero is constructive, whereas normalizing an
arbitrary vector to a unit vector is not: normalization needs a lower
bound for its length.

\subsection{Constructing an integral directly from its finite sums}

Fix a PR rational evaluator \(f(a,x,p)\) on rational points of
\([0,1]^d\), with regular names for its values.  Suppose a nondecreasing
PR modulus \(\mu\) is supplied and its arithmetic correctness is proved:
\begin{equation}\label{eq:uc}
 \|x-y\|_\infty\le2^{-\mu(j)}
 \quad\Longrightarrow\quad |f(x)-f(y)|\le2^{-j}.
\end{equation}
The notation on the right abbreviates finite name comparisons.  Replacing
\(\mu(j)\) by \(j+\max_{i\le j}\mu(i)\) makes it nondecreasing if
necessary.  For simplicity take \(d\ge1\); dimension zero is evaluation.
Set \(m_k=\mu(k+3)\), and define
\begin{equation}\label{eq:integral-name}
 S_k=2^{-dm_k}
 \sum_{0\le i_1,\ldots,i_d<2^{m_k}}
 f\left(a,(i_1,\ldots,i_d)2^{-m_k},k+3\right).
\end{equation}
This is a PR rational computation with exactly \(2^{dm_k}\) summands.

\begin{theorem}[Finite-grid integral]\label{thm:integral}
For \(k\le l\), the sums \eqref{eq:integral-name} satisfy
\[
 |S_k-S_l|\le3\cdot2^{-k-3}<2^{-k}.
\]
Thus they directly define a regular integral name.  The construction and
its correctness are available in \ipra under the stated name and modulus
contracts.
\end{theorem}

\begin{proof}
The finer dyadic grid at level \(m_l\) refines the grid at \(m_k\).
Group its left-endpoint samples according to their coarse cubes.  Each
fine sample and the coarse left endpoint have sup distance at most
\(2^{-m_k}\).  By \eqref{eq:uc} their represented values differ by at
most \(2^{-k-3}\).  The actual rational evaluations introduce errors at
most \(2^{-k-3}\) and \(2^{-l-3}\).  The weights are nonnegative and
sum to one.  Taking their finite weighted sum bounds the difference by
the sum of those three errors, proving the assertion.

To express the argument entirely with rational observations, fix \(k,l\),
the two grids, and a further arbitrary accuracy.  Instantiate the modulus
and name contracts at all the finitely many required points and
precisions, apply the finite weighted-sum inequality, and then use
Lemma~\ref{lem:rational-separation} with the common denominator of
\(S_k,S_l\) and the displayed rational bound.  This removes the
additional arbitrarily fine error and proves the exact non-strict
inequality stated.  The invariant is the
partial weighted-error bound.  No induction on a formula asserting
continuity at all points or convergence at all precisions is used.
\end{proof}

Linearity is proved by applying the finite sum identity on a common
refinement and the same error estimate.  Under ordinary semantics the
name is the Riemann integral: every tagged partition of sufficiently small
mesh is compared to a common refinement using \eqref{eq:uc}, giving the
usual Riemann-sum criterion.  Internally the finite estimate is the
definition and the theorem.  This construction is reusable by the
calculus, Gamma, contour, and geometric packets; each still must provide
the correct modulus and domain bounds for its integrand.

\subsection{Finite positivity and finite cover certificates}

Evaluation at represented points is itself a construction, which must
first be specified.  If \(X(b,p)\) is a fixed PR coordinate-name family
for a point of \([0,1]^d\), put
\[
 \widetilde f(a,b,k)=
 f\bigl(a,\operatorname{clip}_{[0,1]^d}
       X(b,\mu(k+3)+2),k+3\bigr).
\]
Clipping is coordinatewise rational maximum and minimum.  For \(k\le l\),
the two input samples differ by at most
\(2^{-\mu(k+3)-1}\); the modulus and the two evaluation errors
therefore give
\(|\widetilde f_k-\widetilde f_l|\le3\cdot2^{-k-3}\).
This defines the continuous extension as a regular name.  Two equivalent
coordinate names give equivalent output names by the same finer-stage
estimate.  No function sort or choice principle has been introduced.

Given a modulus \eqref{eq:uc}, one can verify positivity throughout a box
by a finite computation.  For a chosen \(j\), take a dyadic grid with
mesh \(h=2^{-\mu(j)-1}\), including its boundary vertices.  If every
grid value sampled at precision \(j\) is at least \(3\cdot2^{-j}\),
then every represented point in the box has value at least \(2^{-j}\).
Indeed, approximate the point to within \(h/4\), clamp its coordinates
to the box, and round to a nearest grid vertex.  Its distance to that
vertex is at most \(3h/4<2^{-\mu(j)}\).  The modulus loses at most
\(2^{-j}\), and the grid evaluation loses at most another \(2^{-j}\).
Everything used to choose the vertex is rational and bounded.

Likewise, a finite list of open sup-norm balls with rational centres and
rational radii can be certified to cover
the cube by checking that every vertex of a mesh-\(h\) grid belongs to
one listed ball with a spare radius greater than \(h\).  Rounding a
represented point as above and choosing the first successful ball
produces a PR cover selector.  Finite search verifies the certificate.

These are useful substitutes for particular compactness steps.  They do
not prove that every qualitatively given open cover has a computable
certificate, or that pointwise positivity supplies a uniform positive
lower bound.  The mesh and the margin are mathematical information that
must be established.  In the branch construction below they are supplied
by explicit coefficient inequalities.
